\chapter{The Albanese universal property of \(\Pic^0_{C/k}\)}
\label{chap:Albanese_AlbaneseUP}
% archon:covers AlgebraicJacobian/Albanese/AlbaneseUP.lean

Let \(k\) be an algebraically closed field, let \(C\) be a smooth proper
geometrically irreducible curve over \(k\), and let \(P_0\in C(k)\).  Write
\(g\) for the genus of \(C\) and assume \(g>0\).  This chapter proves that
the Abel--Jacobi map
\[
  \iota_{P_0}\colon C\longrightarrow \Pic^0_{C/k}
\]
is the Albanese morphism of the pointed curve \((C,P_0)\).  The proof uses
the canonical map from the \(g\)-th symmetric power of \(C\) to its
Jacobian.  In positive characteristic, generic injectivity of this map does
not by itself imply birationality: a purely inseparable map may also have
one geometric point in its generic fibre.  The differential calculation
that excludes this possibility is therefore included explicitly.
Following Milne's Jacobian construction, the universal-property input is
\cite{milne-av}.

\section{The Jacobian and the Abel--Jacobi map}
\label{sec:albaneseup_setup}

The construction of the identity component of the Picard scheme gives
\[
  J:=\Pic^0_{C/k}.
\]
It is a group scheme by \cref{thm:pic0_group_structure}, smooth by
\cref{thm:pic0_smooth}, proper by \cref{thm:pic0_proper}, and geometrically
irreducible by \cref{thm:pic0_geom_irred}.

\begin{definition}
  [The Abel--Jacobi morphism]
  \label{def:abel_jacobi_morphism}
  \dcref{custom}
  \lean{AlgebraicGeometry.Pic0.abelJacobi}
  \uses{def:pic_scheme, thm:pic0_group_structure}

  The relative degree-zero line bundle associated with the rigidified
  diagonal on \(C\times C\) defines a classifying morphism
  \(\iota_{P_0}\colon C\to J\), called the \emph{Abel--Jacobi morphism}.
\end{definition}

\begin{lemma}
  [Formula and basepoint of the Abel--Jacobi morphism]
  \label{lem:abel_jacobi_morphism}
  \dcref{custom}
  \uses{def:abel_jacobi_morphism}

  There is a regular morphism
  \[
    \iota_{P_0}\colon C\longrightarrow J,
    \qquad Q\longmapsto [\mathcal O_C(Q-P_0)],
  \]
  and \(\iota_{P_0}(P_0)=0_J\).
\end{lemma}

\begin{proof}
  On \(C\times C\), viewed over the second factor, consider the
  rigidified relative degree-zero line bundle represented by
  \[
    \mathcal O_{C\times C}
      (\Delta-P_0\times C-C\times P_0).
  \]
  Its restriction over \(Q\in C(k)\), modulo a line bundle pulled back
  from the base, is \(\mathcal O_C(Q-P_0)\).  The moduli property of
  \(\Pic^0_{C/k}\) therefore gives the displayed morphism.  At \(Q=P_0\)
  the fibre is the trivial line bundle, whose class is the identity of
  \(J\).
\end{proof}

\section{Symmetric powers and their canonical maps}
\label{sec:albaneseup_symmetric_powers}

\begin{definition}
  [Symmetric power of a curve]
  \label{def:symmetric_power_curve}
  \dcref{custom}

  For \(r\ge 1\), the \emph{\(r\)-th symmetric power} \(C^{(r)}\) is the
  quotient of \(C^r\) by permutation of the factors, with quotient morphism
  \(\pi_r\colon C^r\to C^{(r)}\).  Every symmetric morphism from \(C^r\) to a
  \(k\)-scheme factors uniquely through \(\pi_r\).
\end{definition}

\begin{lemma}
  [Properties of the symmetric-power projection]
  \label{lem:symmetric_power_projection_properties}
  \dcref{custom}
  \uses{def:symmetric_power_curve}

  The quotient morphism \(\pi_r\colon C^r\to C^{(r)}\) is finite,
  surjective, and separable.
\end{lemma}

\begin{proof}
  On an invariant affine chart the target ring is the ring of
  \(S_r\)-invariants.  Every element of the source ring is integral over this
  invariant ring, so the quotient is finite; the topological quotient
  description gives surjectivity.  Finally, the induced extension of
  function fields is the finite invariant-field extension for a group of
  automorphisms, and is therefore separable.
\end{proof}

\begin{definition}
  [The symmetric-power carrier]
  \label{def:symmetric_power_curve_carrier}
  \dcref{custom}
  \lean{AlgebraicGeometry.Pic0.SymmetricPower}
  \uses{def:symmetric_power_curve}

  The underlying \(k\)-scheme of \(C^{(r)}\) is denoted
  \(\operatorname{SymmetricPower}(C,r)\).
\end{definition}

\begin{lemma}
  [Nonsingularity of a symmetric power of a curve]
  \label{lem:symmetric_power_nonsingular}
  \dcref{custom}
  \uses{def:symmetric_power_curve}

  If \(C\) is nonsingular, then \(C^{(r)}\) is nonsingular.
\end{lemma}

\begin{proof}
  At the image of \((Q,\ldots,Q)\), choose a parameter \(X\) at \(Q\).
  The completed local ring of \(C^r\) is
  \(k[[X_1,\ldots,X_r]]\), and the completed local ring of the quotient is
  its invariant subring.  The fundamental theorem on symmetric functions
  identifies this with
  \(k[[\sigma_1,\ldots,\sigma_r]]\), which is regular.  At a divisor with
  several distinct support points, the completed local ring is the
  completed tensor product of the corresponding invariant power-series
  rings, and is again regular.
\end{proof}

\begin{definition}
  [The symmetric sum morphism]
  \label{def:symmetric_product_av_map}
  \dcref{custom}
  \lean{AlgebraicGeometry.Pic0.symmetricPowerAVMap}
  \uses{def:symmetric_power_curve_carrier}

  For a morphism \(\varphi\colon C\to A\) to an abelian variety, denote the
  induced morphism from the symmetric-power carrier by
  \(\varphi^{(r)}\colon C^{(r)}\to A\).
\end{definition}

\begin{lemma}
  [Symmetrisation of a morphism to an abelian variety]
  \label{lem:symmetric_product_av_map}
  \dcref{custom}
  \uses{def:symmetric_power_curve, def:symmetric_product_av_map}

  Let \(A\) be an abelian variety and let \(\varphi\colon C\to A\) be a
  morphism.  There is a unique morphism
  \(\varphi^{(r)}\colon C^{(r)}\to A\) satisfying
  \[
    \varphi^{(r)}(P_1+\cdots+P_r)
      =\varphi(P_1)+\cdots+\varphi(P_r).
  \]
\end{lemma}

\begin{proof}
  The morphism \(C^r\to A\) obtained by applying \(\varphi\) in every
  coordinate and then adding is invariant under every permutation, because
  the group law of \(A\) is commutative.  The universal property of
  \(C^{(r)}\) gives the unique factorisation.
\end{proof}

\begin{definition}
  [The canonical symmetric-power morphism to the Jacobian]
  \label{def:symmetric_product_to_jacobian_morphism}
  \dcref{custom}
  \lean{AlgebraicGeometry.Pic0.symmetricPowerToJacobian}
  \uses{def:abel_jacobi_morphism, def:symmetric_product_av_map}

  The Lean construction obtained by symmetrising \(\iota_{P_0}\) is denoted
  \[
    f^{(r)}\colon C^{(r)}\longrightarrow J.
  \]
\end{definition}

\begin{lemma}
  [Divisor description of the canonical symmetric-power morphism]
  \label{lem:symmetric_product_to_jacobian_formula}
  \dcref{custom}
  \uses{def:symmetric_product_to_jacobian_morphism,
    lem:abel_jacobi_morphism, lem:symmetric_product_av_map}

  On effective divisors, the canonical morphism is given by
  \[
    f^{(r)}(D)=[\mathcal O_C(D-rP_0)].
  \]
  Denote its reduced closed image by \(W^r\).  The fibre through an
  effective divisor \(D\) is its complete linear system \(|D|\), a projective
  space of dimension \(h^0(D)-1\).
\end{lemma}

\begin{proof}
  Pulling back along \(\pi_r\) gives the sum of the \(r\) Abel--Jacobi
  classes, hence the displayed divisor class.  Two effective divisors have
  the same class precisely when they are linearly equivalent.  The functor
  of effective divisors linearly equivalent to \(D\) is represented by the
  projectivisation of \(H^0(C,\mathcal O_C(D))\), giving the fibre statement.
\end{proof}

\begin{lemma}
  [General divisors of degree at most the genus]
  \label{lem:symmetric_product_general_divisor}
  \dcref{custom}
  \uses{lem:symmetric_product_to_jacobian_formula}

  If \(1\le r\le g\), there is a dense open subset of \(C^{(r)}\) on which
  \(h^0(D)=1\).  Thus the fibres of \(f^{(r)}\) over this open consist of one
  geometric point.
\end{lemma}

\begin{proof}
  \uses{lem:symmetric_power_projection_properties,
    thm:adelic_serreDuality, thm:adelic_riemannRoch}
  If \(h^1(D)>0\), Serre duality identifies
  \[
    H^1(C,\mathcal O_C(D+Q))^\vee
  \]
  with the subspace of \(H^0(C,\Omega_C^1(-D))\) consisting of
  differentials that vanish at \(Q\).  Outside the common zero locus of a
  basis of this latter space,
  adding \(Q\) lowers \(h^1\) by exactly one.  Starting with the zero
  divisor, for which \(h^1=g\), and repeating this argument \(r\) times
  gives a dense open subset of \(C^r\) on which
  \(h^1(P_1+\cdots+P_r)=g-r\).  Riemann--Roch then gives
  \[
    h^0(P_1+\cdots+P_r)
      =r+1-g+(g-r)=1.
  \]
  Intersecting its finitely many translates under \(S_r\) gives a dense
  symmetric open with the same property.  Its image in the finite quotient
  \(C^{(r)}\) is the complement of the image of its closed invariant
  complement, hence is open and dense.
  Since the fibre of \(f^{(r)}\) through \(D\) is \(|D|\), it is a single
  point there.
\end{proof}

\begin{lemma}
  [Invariant one-forms on the Jacobian]
  \label{lem:jacobian_invariant_one_forms}
  \dcref{custom}

  Evaluation at the identity is an isomorphism
  \[
    H^0(J,\Omega_J^1)\xrightarrow{\sim}T_0J^\vee.
  \]
\end{lemma}

\begin{proof}
  \uses{thm:pic0_smooth, thm:pic0_proper, thm:pic0_geom_irred}
  A cotangent vector at the identity extends uniquely to a
  translation-invariant one-form on the smooth group scheme \(J\).
  Conversely, for a global one-form \(\eta\), the coefficients of
  \(t_a^*\eta-\eta\), expressed in an invariant frame, vary regularly with
  \(a\in J\).  They are regular functions on the proper geometrically
  irreducible variety \(J\), hence constant; at \(a=0\) they vanish.  Thus
  every global one-form is invariant, and evaluation at the identity gives
  the asserted inverse pair.
\end{proof}

\begin{lemma}
  [Differential of the Abel--Jacobi map]
  \label{lem:abel_jacobi_differential}
  \dcref{custom}
  \uses{def:abel_jacobi_morphism}

  Under the tangent-space identification
  \(T_0J\simeq H^1(C,\mathcal O_C)\) and Serre duality, the transpose of the
  translated differential of \(\iota_{P_0}\) at \(Q\) is evaluation
  \[
    H^0(C,\Omega_C^1)\longrightarrow T_Q^*C,
    \qquad \omega\longmapsto\omega(Q).
  \]
\end{lemma}

\begin{proof}
  \uses{thm:pic0_tangent_space_iso, thm:adelic_serreDuality}
  Moving \(Q\) to first order changes the class
  \([\mathcal O_C(Q-P_0)]\) by the connecting homomorphism associated with
  the principal-parts sequence at \(Q\).  The Serre-duality transpose of
  this connecting homomorphism is the restriction of a differential to its
  value at \(Q\), which is the displayed evaluation map.
\end{proof}

\begin{lemma}
  [One-forms pulled back by the Abel--Jacobi map]
  \label{lem:abel_jacobi_one_forms}
  \dcref{custom}
  \uses{lem:jacobian_invariant_one_forms,
    lem:abel_jacobi_differential}

  Pullback by \(\iota_{P_0}\) is an isomorphism
  \[
    H^0(J,\Omega_J^1)\xrightarrow{\sim}H^0(C,\Omega_C^1).
  \]
\end{lemma}

\begin{proof}
  By \cref{lem:jacobian_invariant_one_forms}, a form on \(J\) is determined
  by its value at the identity.  Under
  \cref{lem:abel_jacobi_differential}, pulling it back along
  \(\iota_{P_0}\) recovers at every \(Q\) the corresponding canonical
  differential on \(C\).  Under the Serre-duality identification this map
  is the identity, hence an isomorphism.
\end{proof}

\begin{lemma}
  [Dimension of the Jacobian]
  \label{lem:jacobian_dimension}
  \dcref{custom}

  The Jacobian has dimension \(g\).
\end{lemma}

\begin{proof}
  \uses{lem:abel_jacobi_one_forms, thm:pic0_smooth}
  Since \(J\) is smooth, its dimension is the dimension of its cotangent
  space at the identity.  Translation identifies that cotangent space with
  the global one-forms on \(J\).  By
  \cref{lem:abel_jacobi_one_forms}, this space is isomorphic to
  \(H^0(C,\Omega_C^1)\), whose dimension is the genus \(g\).
\end{proof}

\begin{proposition}
  [Differentials on a symmetric power]
  \label{prop:symmetric_product_differentials}
  \dcref{custom}
  \uses{lem:symmetric_product_to_jacobian_formula}

  Pullback by \(f^{(r)}\) is an isomorphism
  \[
    H^0(J,\Omega_J^1)\xrightarrow{\sim}
    H^0(C^{(r)},\Omega_{C^{(r)}}^1).
  \]
  Under the canonical identification of either side with
  \(H^0(C,\Omega_C^1)\), the form corresponding to \(\omega\) vanishes at
  an effective divisor \(D\) if and only if
  \(\operatorname{div}(\omega)\ge D\).
\end{proposition}

\begin{proof}
  \uses{lem:symmetric_power_projection_properties,
    lem:abel_jacobi_one_forms}
  A global one-form on \(C^r\) is a sum of pullbacks from the factors.
  Hence the invariant forms are exactly
  \(\sum_i p_i^*\omega\), with \(\omega\in H^0(C,\Omega_C^1)\).
  The quotient map \(\pi_r\) is separable, so pullback embeds
  \(H^0(C^{(r)},\Omega^1)\) into this invariant subspace.  The composite
  \[
    H^0(J,\Omega_J^1)\longrightarrow H^0(C^{(r)},\Omega^1)
      \longrightarrow H^0(C^r,\Omega^1)^{S_r}
  \]
  is the pullback by the sum of the \(r\) copies of the Abel--Jacobi map.
  By \cref{lem:abel_jacobi_one_forms}, it sends a form to
  \(\sum_i p_i^*\omega\) and is an isomorphism.  Consequently both arrows
  in the display are isomorphisms.

  It remains to identify vanishing at a divisor.  First take \(D=rQ\),
  choose a parameter \(X\) at \(Q\), and write
  \(\omega=(a_0+a_1X+\cdots)dX\).  In the completed local ring of
  \(C^{(r)}\) at \(D\), use the elementary symmetric functions
  \(\sigma_1,\ldots,\sigma_r\) in local parameters
  \(X_1,\ldots,X_r\), and put
  \(\tau_j=\sum_iX_i^j\,dX_i\).  The identities
  \[
    \sigma_m\tau_0-\sigma_{m-1}\tau_1+\cdots+(-1)^m\tau_m
      =d\sigma_{m+1}
      \qquad(0\le m<r)
  \]
  show that \(\tau_0,\ldots,\tau_{r-1}\) form a basis of the cotangent
  module.  The descended form is \(\sum_{j\ge0}a_j\tau_j\), so it
  vanishes at \(D\) exactly when
  \(a_0=\cdots=a_{r-1}=0\), equivalently when
  \(\operatorname{ord}_Q(\omega)\ge r\).  The same calculation in the
  completed local product for
  \(D=\sum_Q n_QQ\) gives the stated criterion.
\end{proof}

\begin{lemma}
  [Tangent sequence of the canonical symmetric-power map]
  \label{lem:symmetric_product_tangent_exact}
  \dcref{custom}
  \uses{lem:symmetric_product_to_jacobian_formula}

  Let \(F=|D|\) be the fibre of \(f^{(r)}\) through an effective divisor
  \(D\), and put \(a=f^{(r)}(D)\).  The sequence
  \[
    0\longrightarrow T_DF\longrightarrow T_DC^{(r)}
      \xrightarrow{\,d f^{(r)}_D\,}T_aJ
  \]
  is exact.  In particular, if \(h^0(D)=1\), then
  \(d f^{(r)}_D\) is injective.
\end{lemma}

\begin{proof}
  \uses{prop:symmetric_product_differentials,
    lem:symmetric_power_nonsingular, thm:adelic_serreDuality,
    thm:adelic_riemannRoch}
  The first arrow is injective because the fibre is a closed subscheme, and
  its image lies in the kernel because \(f^{(r)}\) is constant on the
  fibre.  By \cref{prop:symmetric_product_differentials}, the annihilator
  in \(T_aJ^\vee\) of the image of \(d f^{(r)}_D\) is
  \(H^0(C,\Omega_C^1(-D))\).  Serre duality therefore gives
  \[
    \dim\ker(d f^{(r)}_D)=r-g+h^1(D).
  \]
  The fibre \(F\) is the projective space of effective divisors in the
  complete linear system, so it is smooth of dimension \(h^0(D)-1\).
  Thus \(\dim T_DF=\dim|D|=h^0(D)-1\), and Riemann--Roch says
  \(h^0(D)-1=r-g+h^1(D)\).  The included subspaces have the same
  dimension, proving exactness.  When \(h^0(D)=1\), this dimension is zero.
\end{proof}

\begin{lemma}
  [Birationality of the canonical map]
  \label{lem:symmetric_product_to_jacobian}
  \dcref{custom}
  \uses{lem:symmetric_product_to_jacobian_formula}

  For \(1\le r\le g\), the morphism
  \[
    f^{(r)}\colon C^{(r)}\longrightarrow W^r
  \]
  is birational.  In particular, \(W^g=J\), and
  \(f^{(g)}\colon C^{(g)}\to J\) is a birational surjective morphism.
\end{lemma}

\begin{proof}
  \uses{lem:symmetric_power_projection_properties,
    lem:symmetric_product_general_divisor,
    lem:symmetric_product_tangent_exact, lem:jacobian_dimension,
    thm:pic0_geom_irred}
  The product \(C^r\) has dimension \(r\), and the finite surjective map
  \(C^r\to C^{(r)}\) shows that \(C^{(r)}\) also has dimension \(r\).
  By \cref{lem:symmetric_product_general_divisor}, the map is injective on
  geometric points over a dense open subset.  It is therefore generically
  quasi-finite, so its image \(W^r\) has dimension \(r\).  Thus the induced
  extension of function fields has separable degree one.  In positive characteristic
  this still leaves open the possibility of a nontrivial purely
  inseparable extension.  Choose \(D\) in that open.  Then \(h^0(D)=1\),
  so \cref{lem:symmetric_product_tangent_exact} makes
  \(d f^{(r)}_D\) injective.  Since source and image both have dimension
  \(r\), the differential has full rank, and the function-field extension
  is generically separable.  Its inseparable degree is therefore also one.
  The function fields of \(C^{(r)}\) and \(W^r\) are equal, proving
  birationality.

  For \(r=g\), the image \(W^g\) is a closed irreducible subvariety of the
  irreducible variety \(J\) and has dimension \(g=\dim J\), the latter
  equality being \cref{lem:jacobian_dimension}.  Hence
  \(W^g=J\).  Surjectivity follows because \(C^{(g)}\) is proper and its
  image is exactly \(W^g\).
\end{proof}

\section{Descent and the universal property}
\label{sec:albaneseup_descent}

\begin{lemma}
  [Descent through the birational symmetric power]
  \label{lem:descent_through_birational_sigma}
  \dcref{custom}
  \lean{AlgebraicGeometry.Pic0.descentThroughBirationalSigma}
  \uses{def:symmetric_product_av_map,
    def:symmetric_product_to_jacobian_morphism}

  Let \(A\) be an abelian variety and let
  \(\varphi\colon C\to A\) be pointed at \(P_0\).  There is a unique
  regular morphism \(\psi\colon J\to A\) such that
  \[
    \varphi^{(g)}=\psi\circ f^{(g)}.
  \]
\end{lemma}

\begin{proof}
  \uses{lem:symmetric_product_av_map, lem:symmetric_product_to_jacobian,
    lem:symmetric_power_nonsingular, thm:pic0_smooth,
    thm:rational_map_to_av_extends}
  On dense opens \(U\subseteq C^{(g)}\) and \(V\subseteq J\), the map
  \(f^{(g)}|_U\colon U\to V\) is an isomorphism.  Thus
  \(\varphi^{(g)}|_U\circ(f^{(g)}|_U)^{-1}\) defines a rational map
  \(J\dashrightarrow A\).  The Jacobian is smooth and \(A\) is an abelian
  variety, so \cref{thm:rational_map_to_av_extends} extends this rational
  map uniquely to a regular morphism \(\psi\colon J\to A\).  The two
  morphisms \(\varphi^{(g)}\) and \(\psi\circ f^{(g)}\) agree on the dense
  open \(U\); since \(C^{(g)}\) is nonsingular by
  \cref{lem:symmetric_power_nonsingular}, hence reduced, and \(A\) is
  separated, they agree everywhere.  The same dense-open agreement proves
  uniqueness.
\end{proof}

\begin{definition}
  [The basepoint section of the symmetric power]
  \label{def:basepoint_section_symmetric_power}
  \dcref{custom}
  \uses{def:symmetric_power_curve}

  Define
  \[
    s_{P_0}\colon C\longrightarrow C^{(g)},\qquad
    Q\longmapsto Q+(g-1)P_0.
  \]
\end{definition}

\begin{lemma}
  [The basepoint section followed by the canonical map]
  \label{lem:basepoint_section_to_jacobian}
  \dcref{custom}
  \uses{def:basepoint_section_symmetric_power,
    lem:symmetric_product_to_jacobian_formula,
    def:abel_jacobi_morphism}

  One has \(f^{(g)}\circ s_{P_0}=\iota_{P_0}\).
\end{lemma}

\begin{proof}
  At \(Q\in C\), the left side represents
  \[
    \mathcal O_C(Q+(g-1)P_0-gP_0)=\mathcal O_C(Q-P_0),
  \]
  which is the defining class of \(\iota_{P_0}(Q)\).
\end{proof}

\begin{lemma}
  [The basepoint section followed by a symmetric sum]
  \label{lem:basepoint_section_symmetric_map}
  \dcref{custom}
  \uses{def:basepoint_section_symmetric_power,
    lem:symmetric_product_av_map}

  If \(\varphi(P_0)=0_A\), then
  \(\varphi^{(g)}\circ s_{P_0}=\varphi\).
\end{lemma}

\begin{proof}
  At \(Q\in C\), the left side is
  \(\varphi(Q)+(g-1)\varphi(P_0)=\varphi(Q)\).
\end{proof}

\begin{lemma}
  [Equivalence of the two factorisation equations]
  \label{lem:albanese_eq_iff_symmetricPower_eq}
  \dcref{custom}
  \lean{AlgebraicGeometry.Pic0.albanese_eq_iff_symmetricPower_eq}
  \uses{def:abel_jacobi_morphism, def:symmetric_product_av_map,
    def:symmetric_product_to_jacobian_morphism}

  Let \(\varphi\colon C\to A\) satisfy \(\varphi(P_0)=0_A\).  For every
  regular morphism \(\psi\colon J\to A\),
  \[
    \varphi=\psi\circ\iota_{P_0}
      \quad\Longleftrightarrow\quad
    \varphi^{(g)}=\psi\circ f^{(g)}.
  \]
\end{lemma}

\begin{proof}
  \uses{lem:abel_jacobi_morphism, lem:symmetric_product_av_map,
    lem:symmetric_product_to_jacobian_formula,
    lem:av_regular_map_is_hom, lem:basepoint_section_to_jacobian,
    lem:basepoint_section_symmetric_map}
  Suppose first that \(\varphi=\psi\circ\iota_{P_0}\).  Evaluating at
  \(P_0\) gives \(\psi(0_J)=0_A\), so
  \cref{lem:av_regular_map_is_hom} shows that \(\psi\) is a homomorphism.
  Pulling both sides of the desired symmetric-power equation back to
  \(C^g\), both send \((P_1,\ldots,P_g)\) to
  \(\sum_i\psi(\iota_{P_0}(P_i))\).  The uniqueness clause in the
  universal property of \(C^{(g)}\) gives the equation on \(C^{(g)}\).

  Conversely, compose the symmetric-power equation with
  \(s_{P_0}\).  The two basepoint-section identities
  \cref{lem:basepoint_section_symmetric_map} and
  \cref{lem:basepoint_section_to_jacobian} identify the resulting equation
  with \(\varphi=\psi\circ\iota_{P_0}\).  Hence
  \(\varphi=\psi\circ\iota_{P_0}\).
\end{proof}

\begin{theorem}
  [Albanese factorisation through \(\Pic^0_{C/k}\)]
  \label{thm:albanese_universal_property}
  \dcref{custom}
  \lean{AlgebraicGeometry.Pic0.albanese_universal_property}
  \uses{def:abel_jacobi_morphism}

  For every abelian variety \(A\) and every morphism
  \(\varphi\colon C\to A\) with \(\varphi(P_0)=0_A\), there exists a unique
  regular morphism \(\psi\colon J\to A\) such that
  \[
    \varphi=\psi\circ\iota_{P_0}.
  \]
\end{theorem}

\begin{proof}
  \uses{lem:descent_through_birational_sigma,
    lem:albanese_eq_iff_symmetricPower_eq}
  By \cref{lem:descent_through_birational_sigma}, there is a unique
  \(\psi\colon J\to A\) satisfying
  \(\varphi^{(g)}=\psi\circ f^{(g)}\).  The equivalence in
  \cref{lem:albanese_eq_iff_symmetricPower_eq}, applied to every candidate
  morphism, transports both existence and uniqueness to the equation
  \(\varphi=\psi\circ\iota_{P_0}\).
\end{proof}

\begin{corollary}
  [Albanese universal property]
  \label{cor:albanese_universal_property_homomorphism}
  \dcref{custom}
  \uses{thm:albanese_universal_property}

  In the situation of \cref{thm:albanese_universal_property}, the unique
  factor \(\psi\colon J\to A\) is a homomorphism of group schemes.
  Consequently \(\iota_{P_0}\) is initial among pointed morphisms from
  \(C\) to abelian varieties.
\end{corollary}

\begin{proof}
  \uses{lem:av_regular_map_is_hom, lem:abel_jacobi_morphism}
  Since \(\iota_{P_0}(P_0)=0_J\) and \(\varphi(P_0)=0_A\), the
  factorisation gives \(\psi(0_J)=0_A\).  The pointed rigidity theorem
  \cref{lem:av_regular_map_is_hom} therefore makes \(\psi\) a group
  homomorphism.  Uniqueness among all regular factors implies uniqueness
  among homomorphisms.
\end{proof}
